%%%%

%%%   Самарский государственный технический университет

%%%   Кафедра прикладной математики и информатики

%%%

%%%   Преамбула для статьи в журнал "Вестник Самарского государственного

%%%   технического университета. Серия: «Физико-математические науки»"

%%%   Ориентирована под MikTEx 2.4 for Win32

%%%

%%%   Хотя сам журнал собирается под GNU/Linux на дистрибутиве TeXLive



\documentclass[11pt,a4paper,twoside]{article}



\usepackage[cp1251]{inputenc}

\usepackage[T2A]{fontenc}

\usepackage[english,russian]{babel}

\usepackage{samgtubib}

\usepackage[dvips]{graphicx}

\usepackage[multiple]{footmisc}



\usepackage{samgtu}

\renewcommand{\VestnikYear}{2009} % Год издания Вестника

\renewcommand{\VestnikNumber}{2\,(19)} % Номер Вестника

\renewcommand{\VestnikMonth}{Ноябрь} % Месяц выхода Вестника

\renewcommand{\VestnikData}{01.11.09} % Дата подписания Вестника в печать

\renewcommand{\VestnikUPL}{26,64} % Усл. печ. л.

\renewcommand{\VestnikUIL}{25,73} % Уч.-изд. л.

\renewcommand{\VestnikRegNumber}{479/09} % Регистрационный номер

\renewcommand{\VestnikOrder}{994} % Номер заказа







%

%

%

%\begin{document}



\renewcommand{\baselinestretch}{0.96}



\renewcommand{\firstpage}{155}

\renewcommand{\lastpage}{161}





\udk{530.145}%Квантовая теория

\msc{81Q05, 81Q93}% Closed and approximate solutions to the Schro"dinger, Dirac, Klein-Gordon and other equations of quantum mechanics; Quantum control; Коды MSC см. http://www.ams.org/msc/



\title{Некоторые вопросы динамики и управления открытыми квантовыми системами} 

{Некоторые вопросы динамики и управления открытыми квантовыми системами}

\author{А.\,Н.~Печень}{П\,е\,ч\,е\,н\,ь~А.\,Н.}

\institute{\mian.}

\email{E-mail}{apechen@gmail.com}



\otherinf{\fullauthor{Александр Николаевич Печень}\regalia{к.ф.-м.н.}\position{научный сотрудник}

\dept{отд. математической физики}}



\keywords{открытые квантовые системы, квантовое управление, некогерентное управление}



\workabstract{В работе содержится краткий обзор некоторых вопросов неравновесной динамики и управления открытыми квантовыми системами, таких как динамика открытых квантовых систем в режимах слабой связи и малой плотности, некогерентное управление и~универсальные отображения Крауса.}



\engtitle{Some topics in dynamics and control of open quantum systems}

\engauthor{A.\,N.~Pechen}

{P\,e\,c\,h\,e\,n~A.\,N.}

\enginstitute{\engmian.}

\otherenginfm{\fullauthor{Alexander~N. Pechen} \regalia{Ph.\,D. (Phys. \& Math.)} \position{Research Fellow}

\dept{Dept. of Mathematical Physics}}



\engabstract{This work reviews several topics in the non-equilibrium dynamics and control of open quantum systems, including the dynamics in the weak coupling and low density regimes, incoherent control and universal Kraus maps.}



\engkeywords{open quantum systems, quantum control, incoherent control}



\date{21/XII/2010} % Дата поступления статьи в редакцию

\revisiondate{13/IV/2011} % В окончательном виде



\maketitle



\Section[N]{Введение}

В настоящее время неравновесная динамика открытых квантовых систем является предметом активного изучения~\cite{Breuer,ALV}. Точная динамика, как правило, является сложной для практических вычислений. 

Однако существуют физически важные режимы, в~которых точная динамика аппроксимируется решениями значительно более простых уравнений. 

Такими режимами являются режимы слабой связи и~малой плотности.



Режим слабой связи описывает квантовую систему (атом, молекулу либо наночастицу), слабо взаимодействующую с~окружением (резервуаром). 

Формально предел слабой связи определяется как предел, при котором константа взаимодействия стремится к нулю, $\lambda\to 0$, время стремится к~бесконечности, $t\to+\infty$, притом согласованным образом, так, что $\lambda^2t={\rm const}\ne 0$. Редуцированная (т.\,е. усреднённая по состоянию резервуара) динамика системы в~пределе слабой связи описывается марковским мастер"=уравнением~\cite{Breuer}, а~динамика полной системы "--- квантовыми стохастическими дифференциальными уравнениями (КСДУ)~\cite{ALV}. 

Предел слабой связи, применённый к динамике полной системы, называется стохастическим пределом. 

В~п.~\hyperlink{Sec:WCL}{2} настоящей работы обсуждаются обобщённые квантовые стохастические дифференциальные уравнения с операторами квантового мультипольного шума, описывающие поправки по $\lambda$ к КСДУ в стохастическом пределе слабой связи~\cite{IDAQP2002,PechenMN}.



Режим малой плотности в теории открытых квантовых систем описывает квантовую частицу, взаимодействующую посредством столкновений с разреженным газом малой плотности $n$. 

Формально предел малой плотности определяется как предел при $n\to 0$, $t\to\infty$ так, что $nt={\rm const}\ne 0$. 

Марковские мастер"=уравнения для редуцированной динамики были выведены с~использованием иерархии уравнений Боголюбова"--~Борна"--~Грина"--~Кирквуда"--~Ивона (ББГКИ)~\cite{Dumcke}, ранее применённой в~методе Боголюбова к~выводу кинетических уравнений для классических и~квантовых газов в~режимах слабой связи и~малой плотности~\cite{Bogoliubov2006}. 

В~п.~\hyperlink{Sec:LDL}{3} настоящей работы обсуждается полная динамика системы и~окружения в~пределе малой плотности, которая описывается КСДУ с~квантовым процессом Пуассона~\cite{LDL:QSDE, LDL:QSDE2}.



Тесно связанными с~динамикой являются вопросы управления квантовыми системами~[\citen{Butkovskiy1984}--\citen{Brif2010}]. 

Задачи управления открытыми квантовыми системами возникают в~различных приложениях в~физике и~химии, таких как создание с~использованием лазеров заданных атомных либо молекулярных состояний~\cite{Rice2000}, лазерное охлаждение~\cite{Letokhov2007}, управление химическими реакциями~\cite{Tannor1983, Dantus2004}, управление электроном в~квантовой точке с~окружением из ядерных спинов~\cite{Ruskov2002}, лазерное управление ориентацией молекул~\cite{Orientation}, квантовые вычисления с~использованием смешанных состояний и~неунитарной динамики~\cite{Tarasov}. 

В~п.~\hyperlink{Sec:control}{4} настоящей работы обсуждается использование мастер"=уравнений, возникающих в~пределах слабой связи и~малой плотности, при~изучении некогерентного управления открытыми квантовыми системами~\cite{ice1} и~создании универсальных отображений Крауса~\cite{Wu07}.



\smallskip

\Section{Открытые квантовые системы}

Гильбертово пространство и~гамильтониан открытой квантовой системы, взаимодействующей с~окружением, имеют вид ${\cal H}={\cal H}_{\rm S}\otimes{\cal H}_{\rm E}$ и $H=H_{\rm S}\otimes\mathbb I+\mathbb I\otimes H_{\rm E}+H_{\rm int}=H_0+H_{\rm int}$, где ${\cal H}_{\rm S}$ и ${\cal H}_{\rm E}$ "--- гильбертовы пространства системы и окружения, $H_{\rm S}$ и $H_{\rm E}$ "--- свободные гамильтонианы системы и~окружения, и $H_{\rm int}$ "--- гамильтониан взаимодействия системы с~окружением. 

В~данной работе начальное состояние системы и~окружения предполагается имеющим вид $\sigma=\rho_0\otimes\omega_{\rm E}$, где $\rho_0$ и $\omega_{\rm E}$ "--- начальные состояние системы и~окружения. 

Полный гамильтониан $H$ индуцирует унитарную эволюцию полной системы, которая в~представлении взаимодействия определяется унитарным оператором эволюции $U(t)=e^{itH_0}e^{-itH}$. 

Редуцированная динамика системы определяется как  $\rho(t)={\rm Tr}_{\rm E}[U(t)\sigma U^\dagger(t)]$.



Если система не взаимодействует с~окружением, то её эволюция является унитарной, $\rho(t)=V_t\rho_0 V_t^\dagger$, где унитарный оператор $V_t$ удовлетворяет уравнению Шрёдингера $\dot V_t=-iH_{\rm S}V_t$. 

Описание эволюции замкнутой системы посредством уравнения Шрёдингера называется {\it динамическим}, в~то время как описание её динамики как унитарного отображения "--- {\it кинематическим}. 

Эволюция открытых квантовых систем $\rho(t)=\Phi_t(\rho_0)$ в~кинематическом описании при отсутствии начальных корреляций с~окружением определяется вполне положительными сохраняющими след отображениями $\Phi_t$ (отображениями Крауса)~\cite{Kraus}. Динамическое описание открытых систем использует различные динамические мастер"=уравнения для матрицы плотности~\cite{Breuer}. Как правило, и~редуцированная, и~полная динамика открытой квантовой системы являются сложными для практических вычислений и~не могут быть точно вычислены для реалистичных физических систем. 

Тем не менее существуют физически важные режимы, в~которых как редуцированная, так и~полная динамика приближённо описываются относительно простыми уравнениями. 

Два таких режима "--- режимы слабой связи и~малой плотности обсуждаются в~следующих пунктах.



\smallskip \hypertarget{Sec:WCL}{}

\Section{Режим слабой связи}

Режим слабой связи описывает открытые квантовые системы, слабо взаимодействующие с~окружением. 

Типичным примером таких систем является атом, взаимодействующий с~электромагнитным полем посредством гамильтониана взаимодействия вида $H_{\rm int}=\lambda [Q\otimes a^{+}(f)+Q^\dagger\otimes a^{-}(f)]$, где $\lambda$ "--- малая константа связи и~$a^\pm(f)$ "--- операторы рождения и уничтожения ($f:\mathbb R^3\to\mathbb C$ "--- форм"=фактор, и~$a^{+}(f)=\int f({\bf k})a_{{\bf k}}^{+} d{\bf k}$, где ${\bf k}\in\mathbb R^3$ есть импульс фотона). 

Начальное состояние окружения предполагается гауссовым с~двухточечной корреляционной функцией $\omega_{\rm E}(a_{\bf k}^{+} a_{\bf k}^{-})=n_{\bf k}\delta({\bf k}-{\bf k}')$, где $n_{\bf k}$ есть плотность фотонов с импульсом ${\bf k}$.



Предел малой связи определяется как предел при $\lambda\to 0$ и $t\to+\infty$ такой, что $\lambda^2 t=\tau={\rm const}$. Редуцированная и полная динамика в пределе определяются как $\rho_0(\tau)=\lim_{\lambda\to 0}\rho(\tau/\lambda^2)$ и $U_0(\tau)=\lim_{\lambda\to 0} U(\tau/\lambda^2)$, где последний предел понимается в~смысле специальных корреляционных функций~\cite{ALV} (заметим, что оператор эволюции $U(\tau/\lambda^2)$ зависит от $\lambda$ не только через перерастяжку времени, но и через наличие $\lambda$ в гамильтониане). 

Редуцированная динамика удовлетворяет мастер"=уравнению вида~\cite{Breuer}

\begin{equation}\label{wclme}

\frac{d\rho_0(\tau)}{d\tau}=-i[H_{\rm WCL},\rho_0(\tau)]+{\cal L}_{\rm WCL}(\rho_0(\tau)),

\end{equation}

где действие супероператора ${\cal L}_{\rm WCL}$ имеет общий вид ${\cal L}_{\rm WCL}(\rho)=\sum_i(2L_i\rho L^\dagger_i-L^\dagger_iL_i\rho-\rho L^\dagger_iL_i)$. Конкретная форма ${\cal L}_{\rm WCL}$ определяется гамильтонианом $H$ и состоянием окружения $\omega_{\rm E}$. Полная динамика системы и окружения удовлетворяет КСДУ с квантовым белым шумом~\cite{ALV}, которое в простейшем случае имеет вид

\[

\frac{dU_0(\tau)}{d\tau}=\sqrt{\gamma_0}Qb_0^{+}(\tau)U_0(\tau)-\sqrt{\gamma_0}Q^\dagger U_0(\tau)b_0(\tau)-\gamma_0Q^{+}QU_0(\tau).

\]

Здесь операторы рождения и уничтожения $b_0^\pm(\tau)$ "--- операторы квантового белого шума с~коммутационными соотношениями

\begin{equation}\label{wn}

[b_0^{-}(\tau),b_0^{+}(\tau')]=\delta (\tau-\tau').

\end{equation}



Так как в физических приложениях константа связи $\lambda$ может быть мала, но отлична от нуля, то $U_0(\tau)$ только приблизительно описывает точную динамику $U(t)$. Более точная аппроксимация требует анализа старших членов (поправок $\lambda$ к пределу слабой связи) следующего разложения:

\[

 U(\tau/\lambda^2) \approx U_0(\tau)+\sum\limits_{n=1}^\infty \lambda^n U_n(\tau).

\]

Уравнения для старших членов данного разложения получены в~\cite{IDAQP2002}. 

В~эти уравнения входят операторы квантового мультипольного шума, определённые как операторнозначные обобщённые функции $b^\pm_k(\tau)$ ($k=1, 2, \dots$) с~коммутационными соотношениями 

\begin{equation}\label{mn}

 [b_k(\tau),b_l^{+}(\tau')]=i^k\delta _{k,l} \frac{\partial^k\delta (\tau-\tau')}{\partial\tau^k}.

\end{equation}

Особенностью данных соотношений является наличие производных дельта"=функции. 

В~результате операторы $2^k$-польного квантового шума для нечётных $k$ являются операторами в~псевдогильбертовых пространствах, а~именно в~пространствах Фока с~индефинитной метрикой~\cite{PechenMN}. 

Обобщённые КСДУ, описывающие поправки к~стохастическому пределу, включают квантовый мультипольный шум. 

В~частности, уравнение для первой поправки в~простейшем случае имеет вид~\cite{IDAQP2002}

\begin{multline*}

\frac{dU_1(\tau)}{d\tau}=\sqrt{\gamma_0}Qb_0^{+}(\tau)U_1(\tau)-\sqrt{\gamma_0}Q^\dagger U_1(\tau)b_0(\tau)-\gamma_0Q^\dagger QU_1(\tau)+\\

 +\sqrt{\gamma_1}Qb_1^{+}(\tau)U_0(\tau)-\sqrt{\gamma_1}Q^\dagger U_0(\tau)b_1(\tau).

\end{multline*}

Решением этого уравнения является $$U_1(\tau)=U_0(\tau)\int_0^\tau d\tau'[Q_{\tau'} b_1^{+}(\tau')-Q^\dagger _{\tau'} b_1(\tau')],$$ где $Q_\tau=U_0^\dagger(\tau)(Q\otimes \mathbb I)U_0(\tau)$.



\smallskip\hypertarget{Sec:LDL}{}

\Section{Режим малой плотности} 

Режим малой плотности описывает системы, взаимодействующие с~разреженным газом посредством столкновений. 

Типичный гамильтониан для этого режима имеет вид $$H_{\rm int}=Q\otimes a^{+}(f)a^{-}(g) +Q^\dagger\otimes a^{+}(g)a^{-}(f).$$ Взаимодействие может быть сильным. 

Начальное состояние газа $\omega_{\rm E}$ предполагается гауссовым с~двухточечной корреляционной функцией $\omega_{\rm E}(a_{{\bf k}}^{+} a_{{\bf k}'}^{-})=\varepsilon  n_{\bf k}\delta ({\bf k}-{\bf k}')$, где $\varepsilon \ge 0$ --- малое число и $\varepsilon  n_{\bf k}$ --- плотность частиц газа с импульсом ${\bf k}\in\mathbb R^3$. Предел малой плотности определяется как $\varepsilon \to +0$ и $t\to+\infty$ так, что $\varepsilon  t=\tau={\rm const}$.



Мастер-уравнение для редуцированной матрицы плотности системы в пределе малой плотности имеет общий вид

\begin{equation}\label{ldlme}

 \frac{d\rho_0(\tau)}{d\tau}=-i[H_{\rm LDL},\rho_0(\tau)]+{\cal L}_{\rm LDL}(\rho_0(\tau))

\end{equation}

с супероператором ${\cal L}_{\rm LDL}$ специального вида, зависящим от матрицы рассеяния $S$ для столкновений системы с одной частицей газа~\cite{Dumcke}.



\smallskip

\begin{newthm}{Теорема \cite{LDL:QSDE2}}

Полная динамика системы и окружения в пределе малой плотности удовлетворяет КСДУ

\[

  dU_0(\tau)=dN_\tau(S-\mathbb I)U_0(\tau),

\]

где $dN_\tau$ --- квантовый процесс Пуассона интенсивности $S-\mathbb I$.

\end{newthm}





\pagebreak 

\hypertarget{Sec:control}{}

\Section{Некогерентное управление}

В~данном пункте рассматриваются некоторые вопросы управления открытыми квантовыми системами. 

В~динамическом описании эволюция управляемой открытой квантовой системы в~режимах слабой связи и~малой плотности определяется мастер"=уравнениями вида~(\ref{wclme}) и~(\ref{ldlme}), где $H_{\rm WCL, LDL}^c$ и ${\cal L}_{\rm WCL,LDL}^c$ зависят от управления $c$ как от параметра. 

Существуют два типа управления. 

Когерентное управление, обозначаемое как $c=u(t)$, влияет на гамильтоновы аспекты динамики, т.\,е. определяет $H_{\rm WCL}^u$ и~$H_{\rm LDL}^u$, и стремится сохранить квантовую когерентность в~системе. 

Примером когерентного управления является поле лазера. 

Некогерентное управление $c=n_{\bf k}$ влияет на негамильтоновы аспекты динамики, т.\,е. на ${\cal L}_{\rm WCL}^{n_{\bf k}}$ и~${\cal L}_{\rm LDL}^{n_{\bf k}}$, и,~как правило, стремится разрушить квантовую когерентность~\cite{ice1}. 

Физическим примером некогерентного управления является некогерентное излучение, для которого $n_{\bf k}$ "--- плотность числа фотонов с~импульсом ${\bf k}$. 

В~кинематическом описании эволюция открытой квантовой системы под действием управления $c$ определяется семейством отображений Крауса $\Phi^t_c$, так что $\rho(t)=\Phi^t_c(\rho_0)$. 

Два широких класса задач управления квантовыми системами формулируются как задачи минимизации функционалов видов $J_1[c]={\rm Tr}[\rho(T) O]$ и $J_2[c]=\|\rho(T)-\rho_{\rm target}\|$, где $T>0$ "--- заданное время окончания управления и~$\rho(T)=\Phi^T_c(\rho_0)$.



В работе~\cite{Wu07} были построены отображения Крауса типа <<все-в-одно>> (т.\,н.~универсальные отображения Крауса), переводящие все состояния системы (с~любой матрицей плотности) в~одно и~то же конечное состояние. 

Универсальное отображение Крауса, переводящее все $\rho_0$ в~одно и~то же $\rho_{\rm f}$, обозначим как $\Phi_{\rho_{\rm f}}$. Пусть $c^*$ "--- оптимальное управление, минимизирующее $J_1[c]$ или $J_2[c]$, и пусть $\rho^*=\Phi^T_{c^*}(\rho_0)$. 

Если $\Phi^T_{c^*}=\Phi_{\rho^*}$ "--- универсальное отображение Крауса, то управление $c^*$ будет оптимальным сразу для всех $\rho_0$. 

Данное свойство определяет важность универсальных отображений Крауса. 

Открытым оставался вопрос генерации произвольных универсальных отображений Крауса в~реальных физических системах. 

Для решения этой задачи применительно к  квантовым системам, слабо взаимодействующим с некогерентным излучением (в пределе слабой связи), автором разработана методика, в~которой используется специальная комбинация некогерентного управления $n_{\bf k}$, определяющего ${\cal L}_{\rm WCL}^{n_{\bf k}}$, и когерентного управления $u(t)$, определяющего $H_{\rm WCL}^{u}$. 

При широких предположениях на $H_{\rm WCL}^{u}$ и ${\cal L}_{\rm WCL}^{n_{\bf k}}$ показано, что любое универсальное отображение Крауса может быть приближённо создано определённой комбинацией некогерентного $n^*_{\bf k}$ и когерентного $u^*(t)$ управлений. При этом построение универсального отображения Крауса $\Phi_{\rho_{\rm f}}$, где $\rho_{\rm f}=\sum_i p_i|\phi_i\rangle\langle\phi_i|$, происходит в~два этапа. 

Сначала посредством некоторого некогерентного управления $n^*_{\bf k}$ произвольное $\rho_0$ переводится в~матрицу плотности $\tilde\rho_{\rm f}=\sum_i p_i|i\rangle\langle i|$, где $|i\rangle$ "--- собственные векторы~$H_{\rm S}$. 

Затем с~помощью определённого когерентного управления $u^*(t)$ реализуется унитарная динамика, переводящая $|i\rangle$ в $|\phi_i\rangle$ и,~соответственно, $\tilde\rho_{\rm f}$ в~$\rho_{\rm f}$. 

При этом произвольное начальное состояние $\rho_0$ переводится в~заданное конечное состояние $\rho_{\rm f}$.



\smallskip

\Section[n]{Заключение} Таким образом, в работе рассмотрены некоторые вопросы неравновесной динамики и~управления открытыми квантовыми системами. 

В режиме слабой связи полная динамика системы и~окружения описывается уравнениями с~квантовым мультипольным шумом, в~режиме малой плотности "--- квантовым стохастическим дифференциальным уравнением с~квантовым процессом Пуассона. 

Рассмотрено некогерентное управление открытыми квантовыми системами и~его применение к~генерации универсальных отображений Крауса.



\thanks{Работа выполнена при поддержке РФФИ (проект \No~11--01--00828--a) и гранта президента РФ (НШ-7675.2010.1).}



\begin{thebibliography}{99}



\Bibitem{Breuer}

\book The theory of open quantum systems 

\by Breuer~H.-P.,  Petruccione~F.

\publ Oxford University Press

\publaddr Oxford

\yr 2002 

\totalpages 648 

\rtransl русск. пер.:~

\by Бройер~Х.-П., Петруччионе~Ф.

\book Теория открытых квантовых систем

\publaddr Ижевск

\publ РХД

\yr 2010

\totalpages 824



\Bibitem{ALV}

\by Accardi~L., Lu~Y.\,G., Volovich~I.\,V.

\book Quantum theory and its stochastic limit

\publaddr Berlin

\publ Sprin\-ger-Verlag

\yr 2002

\totalpages 473



\Bibitem{IDAQP2002}

\paper Quantum multipole noise and generalized quantum stochastic equations

\by Pechen~A.\,N., Volovich~I.\,V.

\jour Infin. Dimens. Anal. Quantum Probab. Relat. Top.

\yr 2002

\issue 4

\vol 5

\pages 441--464

\arxiv 	\href{http://arxiv.org/abs/math-ph/0202046}{math-ph/0202046}







\RBibitem{PechenMN}

\by    Печень~А.\,Н.

\paper Об одном асимптотическом разложении в~квантовой теории

\jour Матем. заметки

\yr 2004

\vol 75

\issue 3

\pages 459--461

\transl англ. пер.:~

\by Pechen'~A.\,N.

\paper On an asymptotic expansion in quantum theory

\jour Math. Notes

\yr 2004

\vol 75

\issue 3

\pages 426--429





\Bibitem{Dumcke}

\paper The low density limit for an $N$-level system interacting with a free Bose or Fermi gas

\by    D\"umcke~R.

\jour Comm. Math. Phys.

\yr 1985

\issue 3

\vol 97

\pages 331--359







\RBibitem{Bogoliubov2006}

\by Боголюбов~Н.\,Н.

\book Собрание научных трудов в~двенадцати томах

\yr 2006

\totalpages 804

\publ Наука

\publaddr М.

\bookvol 5

\volinfo Неравновесная статистическая механика, 1939--1980

\misctransl

\by Bogolyubov~N.\,N.

\book Collection of scientific works in twelve volumes

\yr 2006

\totalpages 804

\publ Nauka

\publaddr Moscow

\bookvol 5

\volinfo Nonequilibrium Statistical Mechanics





\Bibitem{LDL:QSDE}

\paper A stochastic golden rule and quantum Langevin equation for the low density limit

\by Accardi~L., Pechen~A.\,N., Volovich~I.\,V.

\jour Infin. Dimens. Anal. Quantum Probab. Relat. Top.

\yr 2003

\vol 6

\issue 3

\pages 431--453

\arxiv 	\href{http://arxiv.org/abs/math-ph/0206032}{math-ph/0206032}







\Bibitem{LDL:QSDE2}

\paper Quantum stochastic equation for a test particle interacting with a dilute Bose gas

\by Pechen~A.\,N.

\jour J.~Math. Phys.

\yr 2004

\vol 45

\issue 1

\pages 400--417

\arxiv 	\href{http://arxiv.org/abs/math-ph/0303020}{math-ph/0303020}





\RBibitem{Butkovskiy1984}

\by Бутковский~А.\,Г., Самойленко~Ю.\,И.

\book Управление квантовомеханическими процессами

\publaddr М.

\publ Наука

\yr 1984

\totalpages 256

\misctransl 

\by Butkovskiy~A.\,G., Samoylenko~Yu.\,I.

\book Control of quantum mecha\-nical processes

\publ Nauka

\publaddr Moscow

\yr 1984

\totalpages 256



\Bibitem{Rice2000}

\by Rice~S.\,A., Zhao~M.

\book Optical control of molecular dynamics

\publaddr New York

\publ Wiley

\yr 2000

\totalpages 437



\Bibitem{Brumer2003}

\by Brumer~P.\,W., Shapiro~M.

\book Principles of the quantum control of molecular processes

\publaddr Hoboken, NJ

\publ Wiley-Interscience

\yr 2003

\totalpages 354



\Bibitem{Letokhov2007}

\by Letokhov~V.\,S.

\book Laser control of atoms and molecules

\publaddr Oxford, New York

\publ Oxford University Press

\yr 2007

\totalpages 310



\Bibitem{Brif2010}

\paper Control of quantum phenomena: past, present and futu\-re

\by Brif~C., Chakrabarti~R., Rabitz~H.

\jour New Journal of Physics

\yr 2010

\issue 7

\vol 12

\papernumber 075008

\totalpages 69

\arxiv \href{http://arxiv.org/abs/0912.5121}{0912.5121} \mbox{[quant-ph]}



\Bibitem{Tannor1983}

\paper Control of selectivity of chemical reaction via control of wave packet evolution

\by Tannor~D\,J., Rice~S.\,A.

\jour J.~Chem. Phys.

\yr 1985

\issue 10

\vol 83

\papernumber 5013

\totalpages 6







\Bibitem{Dantus2004}

\paper Experimental coherent laser control of physicochemical proces\-ses

\by Dantus~M., Lozovoy~V.\,V.

\jour Chem. Rev.

\yr 2004

\issue 4

\vol 104

\pages 1813--1860







\Bibitem{Ruskov2002}

\paper Quantum feedback control of a solid-state qubit

\by Ruskov~R., Korotkov~A.\,N.

\jour Phys. Rev.~B

\yr 2002

\vol 66

\issue 4

\papernumber 041401(R)

\totalpages 4 

\arxiv \href{http://arxiv.org/abs/cond-mat/0204516}{cond-mat/0204516} [cond-mat.mes-hall]



\Bibitem{Orientation}

\paper Laser-assisted control of molecular orientation at high tempe\-ratures

\by Zhdanov~D.\,V., Zadkov~V.\,N.

\jour Phys. Rev. A

\yr 2008

\vol 77

\issue 1

\papernumber 011401(R)

\totalpages 4



\Bibitem{Tarasov}

\paper Quantum computer with mixed states and four-valued logic

\by Tarasov~V.\,E.

\jour J.~Phys.~A

\yr 2002

\issue 25

\vol 35

\pages 5207--5235

\arxiv \href{http://arxiv.org/abs/quant-ph/0312131}{quant-ph/0312131}





\Bibitem{ice1}

\paper Teaching the environment to control quantum systems

\by Pechen~A., Rabitz~H.

\jour Phys. Rev.~A

\yr 2006

\issue 6

\vol 73

\papernumber 062102

\totalpages 6

\arxiv \href{http://arxiv.org/abs/quant-ph/0609097}{quant-ph/0609097}



\Bibitem{Wu07}

\paper Controllability of open quantum systems with Kraus-map dynamics

\by Wu~R., Pechen~A., Brif~C., Rabitz~H.

\jour J.~Phys.~A

\yr 2007

\issue 21

\vol 40

\pages 5681--5693

\arxiv \href{http://arxiv.org/abs/quant-ph/0611215}{quant-ph/0611215}





\Bibitem{Kraus}

\by Kraus~K.

\book States, effects, and operations

\publaddr Berlin

\publ Springer-Verlag

\yr 1983

\totalpages 151

\serial Lecture Notes in Physics

\vol 190





\end{thebibliography}



\noindent

\hskip6.5cm \thedate \par\noindent

\hskip6.5cm\therevdate\par







\makeengtitlem



\renewcommand{\baselinestretch}{0.9}



\newpage



%\end{document}

